\chapter{Álgebra qutrit de color y conexión finita de calibre}
\label{chap:algebra-color-qutrit}
\index{color!qutrit}\index{operadores de Weyl}
\index{grupo especial unitario}\index{acción de Wilson}

El funcional residual \(\kappa_3\) del
\cref{chap:ocupacion-estadistica} suministra tres clases afines; el atlas
posterior realiza su balance \(12+12+12\) en la
\cref{prop:color-atlas}. Una terna de etiquetas no es todavía una
representación. El paso algebraico requiere declarar la carta qutrit, su
orientación y su carácter de fase. Bajo esas hipótesis, la representación y
su álgebra infinitesimal se obtienen exactamente, sin utilizar masas ni datos
cromodinámicos.

\section{Par de Weyl y álgebra de traza \(0\)}
\label{sec:weyl-color-qutrit}

Sea
\[
 V_{\rm col}=\mathbb C^3,
 \qquad
 \omega=e^{2\pi i/3},
\]
con base ordenada \((e_0,e_1,e_2)\). Definimos
\[
 Xe_j=e_{j+1\bmod3},
 \qquad
 Ze_j=\omega^j e_j.
\]
Se cumplen
\[
 X^3=Z^3=I,
 \qquad
 ZX=\omega XZ.
\]

\begin{theorem}[Base de Weyl y álgebra compleja de color]
\label{thm:weyl-sl3-color}
Para
\[
 W_{ab}:=X^aZ^b,
 \qquad
 (a,b)\in(\mathbb Z/3\mathbb Z)^2,
\]
los nueve operadores \(W_{ab}\) forman una base ortogonal de
\(\operatorname{End}_{\mathbb C}(V_{\rm col})\) respecto del producto de
Hilbert--Schmidt. Los ocho operadores no identitarios poseen traza cero y
\[
 \operatorname{span}_{\mathbb C}
 \{W_{ab}:(a,b)\ne(0,0)\}
 =\mathfrak{sl}_3(\mathbb C).
\]
\end{theorem}

\begin{proof}
La relación de Weyl permite reducir cualquier producto a la forma
\(X^aZ^b\). Un cálculo directo da
\[
 \operatorname{tr}(W_{ab}^{\dagger}W_{cd})
 =3\,\delta_{ac}\delta_{bd}.
\]
Por tanto, los nueve operadores son linealmente independientes; como
\(\dim_{\mathbb C}\operatorname{End}_{\mathbb C}(V_{\rm col})=9\), forman una
base. La traza de \(X^aZ^b\) se anula salvo para
\((a,b)=(0,0)\). Los ocho restantes forman así una base del subespacio de
traza cero, cuya dimensión compleja es ocho.
\end{proof}

\begin{corollary}[Forma real compacta]
\label{cor:su3-color-qutrit}
La involución
\[
 \tau(A):=-A^\dagger
\]
preserva \(\mathfrak{sl}_3(\mathbb C)\), y su conjunto de puntos fijos es
\[
 \mathfrak{sl}_3(\mathbb C)^\tau
 =\{A:A^\dagger=-A,\ \operatorname{tr}A=0\}
 =\mathfrak{su}(3).
\]
Esta forma real posee dimensión ocho.
\end{corollary}

\begin{proof}
Se tiene
\[
 W_{ab}^{\dagger}=\omega^{ab}W_{-a,-b}.
\]
Los ocho índices no nulos forman cuatro pares adjuntos. Para
\[
 \mathcal R=\{(1,0),(0,1),(1,1),(1,2)\},
\]
los operadores
\[
 A_u=W_u-W_u^\dagger,
 \qquad
 B_u=i(W_u+W_u^\dagger),
 \qquad u\in\mathcal R,
\]
son antihermíticos, de traza cero y linealmente independientes sobre
\(\mathbb R\). Constituyen, por tanto, una base real de
\(\mathfrak{su}(3)\).
\end{proof}

La aplicación residual
\[
 \kappa_{\rm col}(r,c)
 =\mathbb C e_{\,r-c\bmod3}
\]
asocia las tres clases del atlas con las tres rectas propias de \(Z\). El
operador \(X\) las permuta cíclicamente y \(Z\) registra su fase. Esta
asignación es exacta una vez fijados
\[
 V_{\rm col}=\mathbb C[\mathbb F_3],
 \quad
 \text{la base delta},
 \quad
 j\longmapsto j+1,
 \quad
 j\longmapsto\omega^j.
\]
Los cambios afines admisibles de la carta producen representaciones
unitariamente conjugadas. La construcción es, por ello, natural hasta
conjugación unitaria; no selecciona por sí sola los nombres físicos de los
colores ni una realización QCD.

\section{Conexión y acción sobre un complejo celular finito}
\label{sec:gauge-finito-color}

\begin{construction}[Teoría gauge finita relativa al complejo]
\label{cons:gauge-finito-color}
Sea \(\mathcal K=(V,E,F)\) una celularización finita orientada compatible con
el grafo de rutas. A cada arista orientada \(e\) se asigna
\[
 U_e\in SU(3),
 \qquad
 U_{\bar e}=U_e^{-1}.
\]
Una transformación \(g:V\to SU(3)\) actúa por
\[
 U_e\longmapsto U_e^g
 =g_{t(e)}U_eg_{s(e)}^{-1}.
\]
Para cada cara \(f\), el producto ordenado
\[
 U_f=\prod_{e\subset\partial f}U_e^{\epsilon(f,e)}
\]
define su holonomía. Elegido un vértice base \(v_f\), se tiene
\[
 U_f^g=g_{v_f}U_fg_{v_f}^{-1}.
\]
\end{construction}

\begin{proposition}[Invariancia gauge y transporte covariante]
\label{prop:wilson-color-finito}
Para \(\beta\ge0\), la acción
\[
 S_{\rm col}^{(\mathcal K,\beta)}[U]
 =\frac{\beta}{3}\sum_{f\in F}
 \left(3-\operatorname{Re}\operatorname{tr}U_f\right)
\]
es gauge-invariante y no negativa. Si
\(\psi_v\in V_{\rm col}\) transforma como
\(\psi_v\mapsto g_v\psi_v\), entonces
\[
 \nabla_e\psi:=U_e\psi_{s(e)}-\psi_{t(e)}
\]
satisface
\[
 \nabla_e^g\psi^g=g_{t(e)}\nabla_e\psi.
\]
\end{proposition}

\begin{proof}
La holonomía cambia por conjugación y la traza es cíclica; de ahí la
invariancia. Como cada \(U_f\) es unitario de dimensión tres,
\(\operatorname{Re}\operatorname{tr}U_f\le3\), lo que prueba la no
negatividad. La covariancia de \(\nabla_e\) resulta sustituyendo las dos leyes
de transformación.
\end{proof}

\section{Positividad de núcleos de clase y descomposición armónica}

La conexión finita precedente no implica por sí sola positividad por
reflexión ni un límite de Yang--Mills. Antes de formular tales transferencias
conviene aislar el control clásico que sí está cerrado. Sea \(G\) un grupo
compacto y \(\widehat G\) su dual unitario. Para cada
\(\rho\in\widehat G\), sean \(d_\rho\) su dimensión y
\(c_\rho(a)\ge0\) coeficientes con decaimiento suficiente, respecto del
crecimiento de \(d_\rho\), para que la serie siguiente converja
absolutamente.

\begin{definition}[Núcleo de clase]
\label{pf84:YM-07}
Se define
\[
 K_a(g,h):=
 \sum_{\rho\in\widehat G}
 c_\rho(a)\,
 \Tr\!\bigl(\rho(g)\rho(h)^*\bigr),
 \qquad g,h\in G.
\]
La elección
\[
 c_\rho(t)=d_\rho e^{-tC_2(\rho)},
 \qquad t>0,
\]
proporciona el ejemplo del núcleo de calor. La no negatividad de los
coeficientes no reemplaza la hipótesis de convergencia absoluta.
\end{definition}

\begin{theorem}[Positividad de Peter--Weyl]
\label{pf84:YM-08}
Para \(g_1,\dots,g_n\in G\) y \(z_1,\dots,z_n\in\CC\),
\[
 \sum_{i,j=1}^n
 \overline z_i z_jK_a(g_i,g_j)\ge0.
\]
\end{theorem}

\begin{proof}
La unitariedad de las representaciones y la convergencia absoluta permiten
reordenar la suma:
\[
 \sum_{i,j}\overline z_i z_jK_a(g_i,g_j)
 =
 \sum_{\rho\in\widehat G}
 c_\rho(a)
 \left\lVert
 \sum_j\overline z_j\rho(g_j)
 \right\rVert_{\mathrm{HS}}^2\ge0.
\]
Un coeficiente espectral negativo puede destruir esta positividad de Gram;
por ello el signo forma parte de las hipótesis.
\end{proof}

\begin{lemma}[Control abeliano de Poisson]
\label{pf84:YM-09}
Para \(G=U(1)\), \(\theta\in\RR\) y \(|a|<1\),
\[
 \sum_{n\in\ZZ}a^{|n|}e^{in\theta}
 =\frac{1-a^2}{1-2a\cos\theta+a^2}.
\]
Si se exigen además coeficientes de Fourier no negativos, se restringe
\(a\) a \(0\le a<1\).
\end{lemma}

\begin{proof}
La igualdad resulta de sumar por separado las dos series geométricas de
índices positivos y negativos y añadir el modo \(n=0\). Para \(a<0\), la
identidad escalar sigue siendo válida, pero ya no constituye un ejemplo de
coeficientes Peter--Weyl no negativos.
\end{proof}

\begin{remark}
Los \cref{pf84:YM-07,pf84:YM-08,pf84:YM-09} son definiciones y resultados
clásicos de análisis armónico. Proporcionan un control positivo para
núcleos de clase; no demuestran positividad Osterwalder--Schrader, brecha de
masa, confinamiento ni límite continuo, y no promueven la teoría gauge
finita de este capítulo a cromodinámica cuántica.
\end{remark}

\begin{exactbox}[title={Cierre algebraico del color}]
Los \cref{thm:weyl-sl3-color,cor:su3-color-qutrit} construyen exactamente la
cadena
\[
\text{residuo ternario}\longrightarrow
\text{qutrit}\longrightarrow
\mathfrak{sl}_3(\mathbb C)\longrightarrow\mathfrak{su}(3).
\]
La \cref{cons:gauge-finito-color,prop:wilson-color-finito} prolonga esa
cadena a una conexión y una acción gauge sobre la celularización finita
\(\mathcal K\).  Éste es el objeto demostrado en el capítulo.  La
comparación con la cromodinámica física se efectúa después y no redefine el
qutrit, el balance \(12+12+12\) ni la representación construida.
\end{exactbox}
